\subsubsection{Second-Order Scalar Equations}\label{subsec:second-order-scalar-interior-regularity}

In order to eliminate many technical details, we will give a proof of an interior regularity result only in the case of a second-order scalar equation. We already gave statements of results for higher-order equations and systems in Section~\ref{sec:general-linear-elliptic-problems} above.

\begin{Theorem}{Interior regularity}
\label{thm:second-order-interior-regularity}
Let $L$ be a uniformly elliptic second-order operator of the form\[L(\bm{x},D)u:=\sum_{i,j=1}^n\frac{\partial}{\partial x_i}a_{ij}(\bm{x})\frac{\partial u}{\partial x_j}+\sum_{i=1}^nb_i(\bm{x})\frac{\partial u}{\partial x_i}+c(\bm{x})u,\]with corresponding bilinear form\[B[v,u]:=-\sum_{i,j=1}^n\int_\Omega a_{ij}(\bm{x})u_{x_j}v_{x_i}\,\mathrm{d}\bm{x}+\sum_{i=1}^n\int_\Omega b_i(\bm{x})u_{x_i}v\,\mathrm{d}\bm{x}+\int_\Omega c(\bm{x})uv\,\mathrm{d}\bm{x}.\]Suppose the coefficients satisfy $a_{ij}\in W^{1,\infty}(\Omega)$, $b_i,c\in L^\infty(\Omega)$ and that $f\in L^2(\Omega)$. Let $u\in H_0^1(\Omega)$ be a weak solution of the Dirichlet problem for $L(D,\bm{x})u=f$. Then $u\in H^2(\Omega^{\prime})$ for every $\Omega^{\prime}\subset\subset\Omega$, and
\begin{equation}\label{eq:second-order-interior-regularity-h2-estimate}
\|u\|_{H^2(\Omega^{\prime})}\leq C(\|u\|_{H^1(\Omega)}+\|f\|_{L^2(\Omega)}).
\end{equation}

\end{Theorem}

\begingroup
\renewcommand{\endProof}{\par\nobreak\hfill$\blacksquare$\par\vspace{0.1cm}}
\begin{Proof}
We begin with the identity
\begin{equation}\label{eq:interior-regularity-principal-weak-identity}
\int_\Omega a_{ij}u_{x_j}v_{x_i}\,\mathrm{d}\bm{x}=\int_\Omega gv\,\mathrm{d}\bm{x}
\end{equation}
for all $v\in H_0^1(\Omega)$, where $g\in L^2(\Omega)$ is given by
\begin{equation}\label{eq:interior-regularity-reduced-right-hand-side}
g:=b_iu_{x_i}+cu-f.
\end{equation}
Now suppose $v\in H^1(\Omega)$ and $\operatorname{supp}v\subset\subset\Omega$ and let\[2|h|<\operatorname{dist}(\operatorname{supp}v,\partial\Omega).\]Then we have (for any $k=1,\ldots,n$) $D_k^{-h}v\in H_0^1(\Omega)$, and thus we can use \eqref{eq:interior-regularity-principal-weak-identity} and the ``differencing by parts formula\textquotesingle\textquotesingle\ \eqref{eq:difference-quotient-integration-by-parts} to get\[\begin{aligned}
\int_\Omega D_k^h(a_{ij}u_{x_j})v_{x_i}\,\mathrm{d}\bm{x}
&=-\int_\Omega a_{ij}u_{x_j}D_k^{-h}v_{x_i}\,\mathrm{d}\bm{x}\\
&=-\int_\Omega gD_k^{-h}v\,\mathrm{d}\bm{x}.
\end{aligned}\]We can now use this and the product rule for difference quotients \eqref{eq:difference-quotient-shifted-product-rule} to get
\begin{equation}\label{eq:interior-regularity-shifted-principal-product-identity}
\int_\Omega a_{ij}(\bm{x}+h\bm{e}_k)(D_k^hu_{x_j})v_{x_i}\,\mathrm{d}\bm{x}=-\int_\Omega(D_k^ha_{ij})u_{x_j}v_{x_i}+gD_k^{-h}v\,\mathrm{d}\bm{x}.
\end{equation}
We can now use Lemma~\ref{lem:difference-quotient-weak-derivative-bound} to estimate this by
\begin{equation}\label{eq:interior-regularity-shifted-principal-form-bound}
\left|\int_\Omega a_{ij}(\bm{x}+h\bm{e}_k)(D_k^hu_{x_j})v_{x_i}\,\mathrm{d}\bm{x}\right|\leq C(\|u\|_{1,2}+\|f\|_2)\|\nabla v\|_2.
\end{equation}
Let $\Omega^{\prime}\subset\subset\Omega$. We now choose a cutoff function $\zeta\in\mathcal{D}(\Omega)$ with the following properties.\begin{enumerate}\item $\zeta\equiv1$ on $\Omega^{\prime}$.\item $|\nabla\zeta|<2/d$ where $d=\operatorname{dist}(\Omega^{\prime},\partial\Omega)$.\end{enumerate}We now use the ellipticity condition \eqref{eq:second-order-ellipticity-negative-coefficient-matrix}, elementary inequalities, and the previous estimate (with $v=\zeta^2D_k^hu$) to obtain\[\begin{aligned}
\theta\int_\Omega|\zeta\nabla D_k^hu|^2\,\mathrm{d}\bm{x}
&\leq-\int_\Omega\zeta^2a_{ij}(\bm{x}+h\bm{e}_k)D_k^hu_{x_i}D_k^hu_{x_j}\,\mathrm{d}\bm{x}\\
&=-\int_\Omega a_{ij}(\bm{x}+h\bm{e}_k)D_k^hu_{x_i}\left[(\zeta^2D_k^hu)_{x_j}-2D_k^hu\zeta\zeta_{x_j}\right]\,\mathrm{d}\bm{x}\\
&\leq C\Bigl[(\|u\|_{1,2}+\|f\|_2)(\|\zeta\nabla D_k^hu\|_2+2\|D_k^hu\nabla\zeta\|_2)\\
&\qquad+\sup_{i,j}\|a_{ij}\|_\infty\|\zeta\nabla D_k^hu\|_2\|D_k^hu\nabla\zeta\|_2\Bigr]\\
&\leq C(\epsilon)[(\|u\|_{1,2}+\|f\|_2+\|D_k^hu\nabla\zeta\|_2)^2+\epsilon\|\zeta\nabla D_k^hu\|_2^2]
\end{aligned}\]for any $\epsilon>0$.

Thus, after making an appropriate choice of $\epsilon$ and rearranging, we have\[\begin{aligned}
\|D_k^h\nabla u\|_{L^2(\Omega^{\prime})}&\leq\|\zeta D_k^h\nabla u\|_{L^2(\Omega)}\\
&\leq C(\|u\|_{H^1(\Omega)}+\|f\|_{L^2(\Omega)}+\|D_k^hu\nabla\zeta\|)\\
&\leq C(\|u\|_{H^1(\Omega)}+\|f\|_{L^2(\Omega)}).
\end{aligned}\]Here we have used Lemma~\ref{lem:difference-quotient-weak-derivative-bound} and our pointwise bound on $|\nabla\zeta|$.

Finally, we can use this estimate on the difference quotients of $\nabla u$ and Lemma~\ref{lem:uniform-difference-quotient-converse} to deduce that $\nabla u\in H^1(\Omega^{\prime})$. The estimate \eqref{eq:second-order-interior-regularity-h2-estimate} follows immediately.
\end{Proof}
\endgroup
